% Options for packages loaded elsewhere
\PassOptionsToPackage{unicode}{hyperref}
\PassOptionsToPackage{hyphens}{url}
\documentclass[
]{article}
\usepackage{xcolor}
\usepackage{amsmath,amssymb}
\setcounter{secnumdepth}{-\maxdimen} % remove section numbering
\usepackage{iftex}
\ifPDFTeX
  \usepackage[T1]{fontenc}
  \usepackage[utf8]{inputenc}
  \usepackage{textcomp} % provide euro and other symbols
\else % if luatex or xetex
  \usepackage{unicode-math} % this also loads fontspec
  \defaultfontfeatures{Scale=MatchLowercase}
  \defaultfontfeatures[\rmfamily]{Ligatures=TeX,Scale=1}
\fi
\usepackage{lmodern}
\ifPDFTeX\else
  % xetex/luatex font selection
\fi
% Use upquote if available, for straight quotes in verbatim environments
\IfFileExists{upquote.sty}{\usepackage{upquote}}{}
\IfFileExists{microtype.sty}{% use microtype if available
  \usepackage[]{microtype}
  \UseMicrotypeSet[protrusion]{basicmath} % disable protrusion for tt fonts
}{}
\makeatletter
\@ifundefined{KOMAClassName}{% if non-KOMA class
  \IfFileExists{parskip.sty}{%
    \usepackage{parskip}
  }{% else
    \setlength{\parindent}{0pt}
    \setlength{\parskip}{6pt plus 2pt minus 1pt}}
}{% if KOMA class
  \KOMAoptions{parskip=half}}
\makeatother
\setlength{\emergencystretch}{3em} % prevent overfull lines
\providecommand{\tightlist}{%
  \setlength{\itemsep}{0pt}\setlength{\parskip}{0pt}}
\usepackage{bookmark}
\IfFileExists{xurl.sty}{\usepackage{xurl}}{} % add URL line breaks if available
\urlstyle{same}
\hypersetup{
  pdftitle={El desorejamiento se traduce en ordenacismo. La relación con la vidriería},
  hidelinks,
  pdfcreator={LaTeX via pandoc}}

\title{El desorejamiento se traduce en ordenacismo. La relación con la
vidriería}
\author{}
\date{}

\begin{document}
\maketitle

Primer párrafo.

\textbackslash begin\{flushright\}\\
Segundo párrafo alineado a la derecha.\\
\textbackslash end\{flushright\}

El título de arriba, que parece no tener sentido, en realidad es el
nombre del capítulo cinco de la obra de Balandier (1999) y la relación
que pueda tener con el tema del videojuego. Sin embargo, la forma en la
que está escrito responde a una técnica del movimiento de \emph{oulipo}
que se llama N+7, y que se caracteriza por sustituir cada sustantivo del
texto por el séptimo que le sigue en el diccionario, por lo que en
realidad el nombre original es: "El orden se traduce en desorden. La
relación con el videojuego". Este título original, ordenado y con un
mayor sentido, nos plantea entonces dos preguntas: la primera es qué
relación puede tener el orden y el desorden y qué resulta de dicha
conformación; una segunda pregunta es cómo el videojuego, o qué parte
del videojuego, puede ser entendido a partir de dicha relación. Es
necesario entonces entender primero qué se entiende por orden, desorden
y cómo es visto por una ciencia epistémica y una mitología conformadora
de cultura; esto, posteriormente, nos ayudará a entender cómo, dentro de
diferentes prácticas culturales, se maneja de forma paralela o continua
el orden y el desorden. Después, el ver cómo se manejan estos conceptos
dentro de la socialización de la cultura nos permitirá ahora ver cómo es
utilizado en el videojuego y la práctica del juego en línea; es decir,
cómo se subvierten los conceptos orden/desorden una vez que el jugador
entra dentro del proceso lúdico e interactúa dentro de un mundo virtual,
práctica con más interés e impacto que el saber cómo el perder una oreja
en una vidriería se traduce en un mayor apego al orden y las normas.

De forma temporal la relación orden y desorden se puede entender de dos
formas: como un espacio de posibilidad en la que los dos se dan cabida o
de forma continua, en donde ambos, orden y desorden, son condiciones de
posibilidad mutuas. En el primero entra la caracterización de la
inversión, en donde orden y desorden conviven de manera simultánea;
mientras en el segundo, la conversión, nos habla de cómo al desorden le
sigue el orden.

inversión y cpnversion

En el primero, al igual que

\end{document}
